\documentclass[11pt, a4paper]{article}
\usepackage[utf8]{inputenc}
\usepackage[T1]{fontenc}
\usepackage[brazilian]{babel}
\usepackage{amsmath}
\usepackage{amsfonts}
\usepackage{amssymb}
\usepackage[margin=1in]{geometry}
\usepackage{hyperref}
\hypersetup{
    colorlinks=true,
    linkcolor=blue,
    filecolor=magenta,
    urlcolor=cyan,
}

\title{Fusão Latente de Foundation Models (LLM + Séries Temporais) para Predição Financeira Multimodal}

\author{
Rodrigo Banin Ferraz de Camargo\\
\textit{Instituto de Computação -- UNICAMP}
\and
Allan Mariano de Souza\\
\textit{Professor Doutor I. -- Instituto de Computação -- UNICAMP}
}

\begin{document}

\maketitle

\begin{abstract}
Esta proposta de Iniciação Científica tem como objetivo investigar arquiteturas multimodais para predição financeira baseadas na fusão latente entre Large Language Models (LLMs) e Foundation Models para séries temporais. Enquanto séries temporais de mercado capturam dinâmicas quantitativas contínuas, textos financeiros — como notícias e comunicados sobre ativos — carregam informações semânticas sobre eventos e expectativas que influenciam diretamente o comportamento dos preços. Propõe-se um framework de fusão profunda e bidirecional de representações latentes, utilizando mecanismos de atenção cruzada, visando melhorar desempenho preditivo, explicabilidade e robustez sob drift temporal em cenários reais de mercado.
\end{abstract}

\section{Introdução}

A aplicação de aprendizado de máquina em finanças tem crescido significativamente nas últimas décadas, impulsionada pela disponibilidade de grandes volumes de dados de mercado e pelo avanço de modelos de alta capacidade. Entre os problemas centrais estão a predição de retornos, volatilidade, detecção de eventos anômalos e avaliação de risco. Tradicionalmente, esses problemas são abordados a partir de séries temporais numéricas, como preços, volumes e indicadores técnicos.

Entretanto, o comportamento do mercado financeiro não é determinado apenas por padrões históricos de preços. Informações exógenas expressas em linguagem natural — como notícias financeiras, comunicados corporativos e análises de mercado — exercem influência direta sobre expectativas e decisões dos agentes econômicos. Large Language Models (LLMs) tornaram-se ferramentas poderosas para modelar esse tipo de informação textual, permitindo a extração de representações semânticas ricas.

Apesar disso, a maioria das abordagens existentes trata dados textuais e séries temporais de forma separada ou por meio de fusão superficial. Este projeto propõe uma abordagem multimodal que integra profundamente essas duas fontes de informação no espaço latente, permitindo que sinais textuais e temporais interajam explicitamente durante o processo de predição.

\section{Problema e Lacuna}

Modelos baseados exclusivamente em séries temporais apresentam limitações ao ignorar eventos externos que não se manifestam imediatamente nos preços. Por outro lado, abordagens puramente textuais não capturam a estrutura temporal e os regimes quantitativos do mercado. Técnicas comuns de fusão, como concatenação de features ou ensemble tardio, não permitem modelar dependências condicionais entre texto e dinâmica temporal.

A literatura recente carece de arquiteturas que realizem \textbf{fusão latente profunda e bidirecional} entre representações textuais e séries temporais contínuas, especialmente utilizando modelos fundacionais. Essa lacuna limita a capacidade dos modelos de capturar interações complexas entre narrativa financeira e comportamento quantitativo, além de reduzir a robustez a mudanças de regime e drift temporal.

\section{Objetivos}

\subsection{Objetivo Geral}

Desenvolver e validar um framework multimodal de fusão latente entre Large Language Models e Foundation Models para séries temporais financeiras, aplicado à predição de variáveis de mercado.

\subsection{Objetivos Específicos}

\begin{itemize}
    \item Projetar um módulo de alinhamento e fusão latente para dados textuais e séries temporais.
    \item Avaliar o impacto da fusão profunda no desempenho preditivo em tarefas financeiras.
    \item Investigar mecanismos de explicabilidade baseados em atenção cruzada entre modalidades.
    \item Analisar a robustez do modelo frente a drift temporal e mudanças de regime de mercado.
\end{itemize}

\section{Metodologia}

\subsection{Foundation Models}

\subsubsection{Modelos para Séries Temporais}

Para modelar a dinâmica quantitativa do mercado, serão utilizados Foundation Models para séries temporais, capazes de processar janelas contínuas de dados financeiros, como preços e volumes. Dada uma janela temporal $\mathbf{x}_i^{ts}$, extrai-se uma representação latente:

\[
\mathbf{z}_i^{ts} = f_{TS}(\mathbf{x}_i^{ts}; \theta_{ts}) \in \mathbb{R}^{d_z}
\]

Esses modelos capturam padrões temporais, regimes de mercado e dependências de curto e médio prazo.

\subsubsection{Modelos de Linguagem}

Os dados textuais associados a cada ativo — como notícias financeiras — são processados por um Large Language Model pré-treinado, gerando embeddings semânticos:

\[
\mathbf{z}_i^{llm} = f_{LLM}(\mathbf{x}_i^{text}; \theta_{llm}) \in \mathbb{R}^{d_h}
\]

Essas representações capturam informações sobre eventos, sentimento e contexto econômico.

\subsection{Módulo de Alinhamento e Fusão Latente}

As representações latentes são projetadas para um espaço comum:

\[
\tilde{\mathbf{z}}_i^{ts} = W_{ts}\mathbf{z}_i^{ts} + \mathbf{b}_{ts}, \quad
\tilde{\mathbf{z}}_i^{llm} = W_{llm}\mathbf{z}_i^{llm} + \mathbf{b}_{llm}
\]

Em seguida, aplica-se um mecanismo de atenção cruzada bidirecional, permitindo que informações textuais influenciem a interpretação da dinâmica temporal e vice-versa.

A representação fundida final é dada por:

\[
\mathbf{z}_i^{fused} =
\text{LayerNorm}\left(
\alpha(\tilde{\mathbf{z}}_i^{ts} + \mathbf{z}_i^{ts \to llm}) +
(1-\alpha)(\tilde{\mathbf{z}}_i^{llm} + \mathbf{z}_i^{llm \to ts})
\right)
\]

Uma regularização de alinhamento é adicionada para incentivar coerência semântica entre modalidades:

\[
\mathcal{L}_{align} =
\frac{1}{N}\sum_{i=1}^{N}
\|\tilde{\mathbf{z}}_i^{ts} - \tilde{\mathbf{z}}_i^{llm}\|_2^2
\]

\subsection{Cabeça de Predição e Explicabilidade}

A tarefa principal considerada será a \textbf{predição de variáveis financeiras futuras}, como volatilidade ou direção de retorno em um horizonte $T$. A predição é dada por:

\[
\hat{y}_i = W_2 \text{ReLU}(W_1 \mathbf{z}_i^{fused} + \mathbf{b}_1) + b_2
\]

A função de perda utilizada pode ser erro quadrático médio (regressão) ou perda logística (classificação direcional).

A explicabilidade será investigada por meio da análise dos pesos de atenção cruzada, attribution maps temporais e métodos pós-hoc como SHAP, permitindo identificar quais eventos textuais e segmentos temporais mais influenciam a predição.

\subsection{Treinamento e Robustez}

O treinamento ocorre em duas fases: (i) alinhamento multimodal e (ii) treinamento end-to-end. A perda total é:

\[
\mathcal{L}_{total} = \mathcal{L}_{pred} + \lambda_{align}\mathcal{L}_{align}
\]

A robustez será avaliada sob diferentes regimes de mercado, utilizando janelas temporais deslizantes para análise de drift.

\section{Conjunto de Dados}

Serão utilizados datasets públicos multimodais amplamente reconhecidos na literatura:

\begin{itemize}
    \item \textbf{FNSPID}: integra milhões de notícias financeiras com séries temporais históricas de preços de ações.
    \item \textbf{FinMultiTime}: dataset multimodal que alinha texto, tabelas financeiras, séries temporais e indicadores técnicos.
\end{itemize}

Esses datasets permitem avaliar a proposta em cenários realistas de mercado, com grande diversidade temporal e textual.

\section{Referências}

\begin{itemize}
    \item Z. Dong et al. \textit{FNSPID: Financial News and Stock Price Integration Dataset}. GitHub, 2024. Disponível em: \url{https://github.com/Zdong104/FNSPID_Financial_News_Dataset}
    \item W. Wenyan et al. \textit{Multimodal Dataset for Financial Time-Series Forecasting}. HuggingFace, 2024. Disponível em: \url{https://huggingface.co/datasets/Wenyan0110/Multimodal-Dataset-Image_Text_Table_TimeSeries-for-Financial-Time-Series-Forecasting}
    \item W. Wenyan et al. \textit{FinMultiTime: A Four-Modal Dataset for Financial Time-Series Analysis}. arXiv preprint arXiv:2506.05019, 2025. Disponível em: \url{https://arxiv.org/abs/2506.05019}
\end{itemize}

\end{document}
